% Esquisse d'un programme (cote 119, copy 1): archive cote 119, site page 52 (PDF page 53), typescript page 50.
% typescript. Draft compared: EsquisseFr.pdf (text layer):1798-1844.
% AI-transcribed by Claude Opus 5.5 (claude-opus-5-5, adaptive thinking, effort high), Message Batches API, batch esq1 (msgbatch_01LNsogFyF3y6KSp9UbtJtBk), 2026-11-12; not yet checked by eye.
\scanpage{119/52}
\mspage{50}

ceux parmi eux qui ont persévéré suffisamment avec moi, réputé
plus exigeant que d'autres, pour aboutir à un premier travail
noir sur blanc acceptable (certaines fois aussi à un travail mieux
qu'acceptable et plus qu'en seul travail, fait avec goût et jus-%
qu'au bout). Vu la conjoncture, même parmi les rares qui ont per-%
sévéré, plus rares encore seront ceux qui auront l'occasion
d'exercer ce métier, et par là, tout en gagnant leur pain, de
l'approfondir.

10. Depuis l'an dernier, je sens qu'au cours de mon acti-%
vité d'enseignant universitaire, j'ai appris tout ce que j'avais
à en apprendre et enseigné tout ce que je peux y enseigner, et
qu'elle a cessé d'être vraiment utile, à moi-même comme aux au-%
tres. M'obstiner sous ces conditions à la poursuivre encore me
paraîtrait un gaspillage, tant de ressources humaines que de deniers
publics. C'est pourquoi j'ai demandé mon détachement au CNRS (que
j'avais quitté en 1959 comme directeur de recherches frais émoulu,
pour entrer à l'IHES). Je sais d'ailleurs que la situation de
l'emploi est serrée au CNRS comme ailleurs, que l'issue de ma de-%
mande est douteuse, et que si un poste m'y était attribué, ce serait
aux dépens d'un chercheur plus jeune qui resterait sans poste.
Mais il est vrai aussi que cela libérerait mon poste à l'USTL au
bénéfice d'un autre. C'est pourquoi je n'ai pas de scrupule à
faire cette demande, et s'il le faut à revenir à la charge si elle
n'est pas acceptée cette année.

En tout état de cause, cette demande aura été pour moi
l'occasion d'écrire cette esquisse de programme, qui autrement sans
doute n'aurait jamais vu le jour. J'ai essayé d'être bref sans être
sibyllin et aussi, après coup, d'en faciliter la lecture et de
la rendre plus attrayante, en y adjoignant un sommaire. Si malgré
cela elle peut paraître longue pour la circonstance, je m'en excuse.
Elle me paraît courte pour son contenu, sachant bien que dix ans de
travail ne seraient pas de trop pour aller jusqu'au bout du moindre

\margin{52}
